\chapter{Background}
\label{sec:background}

\section{Answer Set Programming}
\label{sec:bg-asp}

A \emph{normal rule} in ASP has the form
\[
a \;\leftarrow\; b_1, \ldots, b_m,\; \mathit{not}\; c_1, \ldots, \mathit{not}\; c_n,
\]
where $a$, the $b_i$ and the $c_j$ are atoms: the head $a$ is derived whenever every positive body atom $b_i$ is derivable and no negative body atom $c_j$ holds. The negation $\mathit{not}$ is \emph{default negation}: $\mathit{not}\; c$ holds in the absence of a derivation for $c$.
Because a rule read this way can justify itself, the meaning of a program is not given by the rules alone but through the notion of a stable model. Fix a set of atoms $X$, the candidate solution, and a ground program $P$. The \emph{reduct} $P^X$ is obtained by deleting every rule whose negative body intersects $X$, and by deleting the negative literals from the rules that survive. The reduct is a positive program, so it has a unique least model, and $X$ is a \emph{stable model} (or \emph{answer set}) of $P$ exactly when it coincides with that least model \cite{gelfond1988stable}.
On the two-rule program
\[
\left\{
\begin{aligned}
a &\leftarrow \mathit{not}\; b\\
b &\leftarrow \mathit{not}\; a
\end{aligned}
\right\}
\]
the candidate $X = \{a\}$ gives the reduct $\{a \leftarrow\}$, whose least model is $\{a\}$: stable.
The candidate $\{a,b\}$ deletes both rules, leaving the least model $\emptyset$: not stable. 
A stable model is then a set of atoms both closed under the rules and justified by them, in which every atom has a non-circular derivation that survives the model's own assumptions about what is false.

In the \emph{guess-and-check} idiom a choice construct opens the solution space, integrity constraints prune the candidates that violate the specification, and the answer sets of the program correspond one-to-one to the solutions of the problem. ASP syntax extends normal rules with first-order variables, arithmetic, conditional literals, and three constructs this thesis uses throughout. 
A \emph{choice rule} such as \texttt{\{ b(X) \} :- x(X).} permits any subset of its head atoms, introducing alternatives that the solver may resolve through propagation or decisions.
An \emph{integrity constraint} \texttt{:- Body.} is a rule with empty head and forbids every model that satisfies $Body$. 
An \emph{aggregate atom} such as \texttt{\#sum}, \texttt{\#count}, \texttt{\#min} or \texttt{\#max} condenses a set of weighted contributions into a single comparable value \cite{clingoGuide}. Aggregates are central to this work: the heuristics under study rank decisions by expressions such as \enquote{the current sum of the elements already placed in a set}, and how and \emph{when} such an aggregate is evaluated is exactly what differentiates the systems compared here.

\section{Ground-and-Solve: gringo, clasp, and Conflict-Driven Search}
\label{sec:bg-cdcl}

clingo \cite{clingoGuide} couples the grounder gringo with the solver clasp into a single system.
The grounder instantiates the first-order program over its Herbrand universe, producing a variable-free program. To mitigate combinatorial explosion, gringo restricts this expansion to rule instances whose positive body is potentially derivable; as a result, the ground program is generally much smaller than a full instantiation, though it must still be materialized in full before search starts. 
The solver clasp \cite{gebser2007clasp} then searches for the stable models of the ground program using \emph{conflict-driven nogood learning} (CDNL), an algorithm that learns from conflict-inducing choices.
The solver maintains a \emph{partial assignment}, a consistent set of literals over the program atoms; it alternates \emph{unit propagation} steps with \emph{decision} steps, and when propagation derives a contradiction, it analyzes the conflict, learns a nogood built around the \emph{First Unique Implication Point} (1UIP), and \emph{backjumps} to the decision level at which that nogood becomes unit. Which atom to decide on, and with which polarity, is delegated to a \emph{decision heuristic}; clasp's default is activity-based, in the family introduced by Moskewicz et al.~\cite{moskewicz2001chaff}, and favors atoms that appear in recent conflicts.

In the single-shot setting considered here, the ground input program is materialized before search begins; additional nogoods are learned during solving. The \emph{grounding bottleneck} is a well-known limitation of ground-and-solve systems \cite{comploi2023domain}; Chapter~\ref{sec:results} measures one concrete manifestation of it, namely heuristic directives whose weight depends on an aggregate value.

\section{Domain-Specific Heuristics in clingo}
\label{sec:bg-heuristics}

clasp allows the modeller to steer its decision heuristic using domain-specific knowledge through the \texttt{\#heuristic} directive \cite{gebser2013heuristics}:

\begin{minted}{prolog}
#heuristic A : Body. [W@P, M]
\end{minted}

Here, $A$ is the target domain atom whose decision choice is steered according to the modifier $M$ (e.g., setting its sign or priority level). The heuristic applies only if the rule $Body$ is satisfied. If multiple heuristic directives target the same atom $A$, clingo selects the one with the highest priority $P$ to determine the atom's modifier and weight. Across distinct target atoms, clasp then prioritizes candidates with higher weight $W$. Alpha reads the same two fields the other way round, as the next section describes.

\section{Lazy Grounding, Alpha, and Heuristics over Partial Assignments}
\label{sec:bg-alpha}

Lazy-grounding solvers invert the order of the two phases. Rather than instantiating the program and then searching it, they instantiate a rule only when the partial assignment reached so far makes it relevant, avoiding the need to materialize the full ground program before search. Exploring unsuccessful branches can, however, require additional rule instantiation and nogood generation. Alpha \cite{weinzierl2017blending} is the external lazy-grounding reference used in this work.

Comploi-Taupe et al.\ extended Alpha with declaratively specified domain-specific heuristics evaluated against the \emph{partial} assignment maintained during solving \cite{comploi2023domain}, and subsequently with aggregates evaluated over that same partial assignment \cite{comploi2023dynamicAggregates}. Because a heuristic is read while the assignment is still incomplete, its semantics differs from clingo's on the following points:

\begin{itemize}
    \item \textbf{Default negation.} \texttt{not a} holds as soon as $a$ is not assigned true, rather than only once $a$ has been assigned false. A guard over an atom the search has not yet reached is therefore open in Alpha and closed in clingo.
    \item \textbf{Dynamic aggregates.} An aggregate is evaluated over its currently true source atoms. For example, if only \texttt{c(2)} and \texttt{c(5)} are true, \texttt{S = \#sum\{Y : c(Y)\}} yields $S=7$, even while other \texttt{c/1} atoms remain unassigned. The heuristic can use this partial sum immediately and update its weight as the assignment changes. In clingo, the same equality binding is grounded over candidate values, and its satisfaction during search follows the ordinary aggregate semantics; it does not simply bind $S$ to the sum of the currently true atoms.
    \item \textbf{Rank.} The annotation $[W@P]$ is an absolute rank over all heuristic instances: the level $P$ decides first, always preferring higher values, and the weight $W$ orders only the instances that share a level. In clingo, as Section~\ref{sec:bg-heuristics} recorded, $P$ arbitrates between directives on the \emph{same} target atom and $W$ alone orders distinct atoms.
\end{itemize}

An encoding written for one system therefore cannot be handed to the other verbatim. Chapter~\ref{sec:methodology} describes how these differences are handled within the two axes of the experimental design: heuristic representation and semantics.

\section{Extending clingo: Propagators and the Heuristic Interface}
\label{sec:bg-propagator}

Standard clingo evaluates grounded heuristic conditions during search, but does not evaluate the symbolic aggregate templates over the currently true atoms in the way required here. It also provides an extension mechanism: \emph{theory propagators}, user-defined components registered with the solver that participate directly in the search loop \cite{gebser2016theorySolving}. In this work, that interface computes heuristic guidance from the current solver trail.

A propagator has four callbacks. During \texttt{init}, it inspects the symbolic atoms of the ground program, maps the atoms it cares about to solver literals, and registers \emph{watches} on them. During search, \texttt{propagate} is invoked when a watched literal becomes true, \texttt{undo} when backtracking retracts it, and \texttt{check} according to the configured check mode, by default on total assignments. Watches trigger propagation callbacks; checks and heuristic decisions have their own invocation points.

The \texttt{clingo::Heuristic} interface inherits from Propagator and extends it with a \texttt{decide} callback, invoked at each decision step \emph{before} the solver's built-in heuristic. The Heuristic object receives the current assignment of the solver and either returns a signed literal, which becomes the next decision, or returns $0$, in which case clasp falls back to its default activity-based heuristic (VSIDS \cite{moskewicz2001chaff}).
This is the mechanism the rest of the thesis builds on. It keeps the declarative heuristic descriptions in memory, evaluates them incrementally against the trail, and answers each \texttt{decide} call with the best applicable target.
